\chapter{Human Papillomavirus (HPV)}
\index{Human Papillomavirus (HPV)}

\section*{Definition and Overview}
Human papillomavirus (HPV) is the name for a group of more than 200 related viruses, some of which are sexually transmitted and infect the skin and mucous membranes of the genital area, mouth, and throat. HPV is the most common sexually transmitted infection in the world. Most people clear the virus without ever knowing they have it, but persistent infection with high-risk types can cause cervical and other cancers, while low-risk types cause genital warts. Vaccination and screening have made HPV-related cancers largely preventable.

\section*{Epidemiology}
Around 80\% of sexually active people acquire HPV at some point, usually in their teens or twenties. At any time, roughly one in ten women has an active genital HPV infection. In many countries, the school-based HPV vaccination program has cut infection with the vaccine types by more than 90\% in young people, and cervical pre-cancer has fallen sharply. some countries are on track to virtually eliminate cervical cancer by 2035. Genital warts, once very common, have nearly disappeared in vaccinated cohorts.

\section*{Aetiology and Pathophysiology}
\begin{itemize}
    \item \textbf{Transmission:} skin-to-skin contact during vaginal, anal, or oral sex; condoms reduce but do not eliminate spread
    \item \textbf{Low-risk types:} HPV 6 and 11 cause about 90\% of genital warts
    \item \textbf{High-risk types:} around 15 types, led by HPV 16 and 18, can cause cancer
    \item \textbf{Pathophysiology:} the virus infects basal cells of the skin or mucosa; most infections are cleared by the immune system within two years
    \item \textbf{Persistence:} if high-risk HPV persists, viral proteins can drive pre-cancerous changes
    \item \textbf{Cancers:} cervical, anal, penile, vaginal, vulval, and oropharyngeal cancers
    \item \textbf{No cure:} management targets warts and cell changes rather than the virus itself
\end{itemize}

\section*{Clinical Features}
\begin{itemize}
    \item Most infections cause no symptoms
    \item Genital warts: small, flesh-coloured bumps on the genitals, anus, or surrounding skin, which may itch or bleed
    \item Cell changes are usually silent and detected by cervical screening
    \item Advanced cervical changes can cause abnormal bleeding or discharge
    \item Throat and anal infections are usually silent
    \item Symptoms of HPV-related cancer depend on the site
    \item Most people never know they have been infected
\end{itemize}

\section*{Differential Diagnosis}
\begin{itemize}
    \item Genital warts are distinguished from skin tags, molluscum contagiosum, and other lesions
    \item Cervical changes are graded by screening and biopsy
    \item Other sexually transmitted infections should be considered
    \item Persistent infection is identified by testing, not symptoms
    \item Most abnormal screening results resolve without treatment
\end{itemize}

\section*{When to Seek Medical Care}
Women attend cervical screening every 5 years from age 25 to 74, and anyone with genital warts, unusual bleeding, or STI concerns should be assessed. HPV vaccination is recommended for adolescents.

\textbf{Contact your local emergency services for:}
\begin{itemize}
    \item Heavy or persistent abnormal vaginal bleeding
    \item Severe pelvic or abdominal pain
    \item Bleeding with collapse or faintness
    \item Any emergency symptoms; HPV itself does not cause acute emergencies
\end{itemize}

\section*{Diagnosis and Investigations}
\begin{itemize}
    \item \textbf{Cervical screening:} HPV testing with cell review as the national screening method
    \item \textbf{Colposcopy:} detailed examination of the cervix after abnormal results
    \item \textbf{Biopsy:} confirms the grade of cell changes
    \item \textbf{Clinical examination:} for genital warts
    \item \textbf{Anal and oral screening:} for at-risk groups
    \item \textbf{HPV typing:} distinguishes high- and low-risk types
\end{itemize}

\section*{Management}
\begin{itemize}
    \item \textbf{Observation:} most infections and low-grade changes clear without treatment
    \item \textbf{Wart treatment:} topical creams, freezing, or laser
    \item \textbf{Monitoring:} repeat screening at recommended intervals
    \item \textbf{Colposcopy and treatment:} loop excision or other procedures for higher-grade cervical changes
    \item \textbf{Cancer treatment:} surgery, radiotherapy, and chemotherapy
    \item \textbf{Vaccination:} prevents infection with the vaccine types
    \item \textbf{Partner discussion:} informing partners is appropriate, though most have been exposed
\end{itemize}

\section*{Complications}
\begin{itemize}
    \item Genital warts causing discomfort and distress
    \item Persistent infection progressing to pre-cancer
    \item Cervical cancer and other HPV-related cancers
    \item Recurrence of warts after treatment
    \item Anxiety and stigma
    \item Fertility and pregnancy considerations after treatment for cell changes
\end{itemize}

\section*{Prognosis and Outcomes}
The prognosis for HPV infection is excellent for the vast majority of people: most clear the virus without consequences, and persistent changes are detected and treated through screening. Vaccinated cohorts have seen near-elimination of warts and cervical pre-cancer, and cervical cancer is projected to become a rare disease in many countries. With vaccination, screening, and treatment, HPV-related cancer is now largely preventable.

\section*{Prevention and Self-Care}
\begin{itemize}
    \item HPV vaccination for girls and boys, ideally before sexual activity
    \item Cervical screening every 5 years from age 25 to 74
    \item Use condoms to reduce transmission
    \item Limit the number of sexual partners
    \item Do not smoke, as smoking increases cancer risk with persistent infection
    \item Seek review of genital bumps or abnormal bleeding
    \item Follow up abnormal screening results promptly
\end{itemize}

\section*{Sources and Further Reading}
The following reputable sources provide further information for healthcare students and patients seeking reliable, current advice about HPV.

\begin{itemize}
    \item Healthdirect. \textit{Human papillomavirus (HPV).} https://www.healthdirect.gov.au/hpv
    \item Cancer Council. \textit{HPV and cervical cancer prevention.} https://www.cancer.org.au/
    \item Australian Government Department of Health. \textit{HPV vaccination and screening.} https://www.health.gov.au/
    \item Family Planning NSW. \textit{HPV information for patients.} https://www.fpnsw.org.au/
    \item World Health Organization. \textit{Human papillomavirus fact sheet.} https://www.who.int/
    \item Centers for Disease Control and Prevention. \textit{HPV infection.} https://www.cdc.gov/
\end{itemize}
